---

### Title
**Hierarchical Optimization Frameworks for Improved Fine-Tuning and Adaptation in Large Language Models**

---

### Abstract
Large Language Models (LLMs) have demonstrated remarkable capabilities across diverse tasks, yet challenges persist in the areas of efficient fine-tuning, robustness in low-resource settings, and generalization to complex domains. Building on recent advancements, this study integrates insights from hierarchical optimization, prototype-based retrieval, and entropy-guided frameworks to propose a novel methodology: **Hierarchical Adaptation for Fine-Tuning and Retrieval (HAFTAR)**. HAFTAR combines first- and zeroth-order optimization, prototype-based task retrieval, and adaptive entropy-guided alignment to achieve optimized performance in multi-task, low-resource, and multilingual settings. Experimental evaluation on benchmarks such as ReasonTabQA, CLaS-Bench, and synthetic reasoning tasks demonstrates significant improvements in task-specific accuracy, semantic relevance, and overall efficiency. The findings highlight the potential of hierarchical hybrid approaches to enhance fine-tuning, task retrieval, and generalization for LLMs.

---

### Introduction
Large Language Models (LLMs) have unlocked unprecedented capabilities in natural language understanding and generation. However, their deployment at scale often encounters challenges, including inefficiencies in fine-tuning, sensitivity to task-specific data scarcity, and poor generalization in multilingual or low-resource contexts. Existing methods, such as supervised fine-tuning (SFT) and reinforcement learning frameworks, often lead to optimization mismatches or performance degradation when dealing with sparse, noisy, or complex datasets. Thus, there is a pressing need for frameworks that can adaptively optimize LLMs while maintaining task specificity, semantic relevance, and computational efficiency.

In this paper, we propose **Hierarchical Adaptation for Fine-Tuning and Retrieval (HAFTAR)**, a unified framework that synergizes hierarchical optimization, task prototype-based retrieval, and entropy-guided alignment. HAFTAR builds upon recent innovations such as Gibbs Initialization with Finite Temperature (GIFT) (Zhao et al., 2026), hierarchical zeroth- and first-order optimization (Hi-ZFO) (Jin et al., 2026), and prototype-based knowledge retrieval (Oh et al., 2026). By leveraging task-specific prototypes and integrating fine-grained zeroth-order corrections, HAFTAR achieves robust performance across diverse tasks and domains.

---

### Related Work
#### Fine-Tuning and Optimization
Recent studies, such as Zhao et al. (2026), highlight the limitations of traditional fine-tuning methods, particularly in their tendency to collapse into suboptimal local minima. Hierarchical frameworks like Hi-ZFO (Jin et al., 2026) have demonstrated the potential of hybrid first- and zeroth-order optimization methods for escaping such minima and improving fine-tuning efficiency.

#### Retrieval and Task Adaptation
Prototype-based knowledge retrieval frameworks (Oh et al., 2026) have emerged as promising solutions for multi-task learning under partial supervision. By embedding task-specific characteristics into prototypes, these frameworks enable robust task adaptation without reliance on fully annotated datasets.

#### Language Model Steering and Alignment
Multilingual benchmarks such as CLaS-Bench (Gurgurov et al., 2026) have underscored the importance of steering LLMs for cross-lingual alignment. Techniques like entropy-guided prompt weighting (Khoury et al., 2026) and Maximum Mean Discrepancy (MMD) guidance (Sani et al., 2026) have shown promise in aligning model outputs with target distributions.

---

### Methods/Experiment
#### Proposed Framework: HAFTAR
HAFTAR integrates three core components for fine-tuning and task adaptation:
1. **Hierarchical Optimization**:
   - Combines first-order (FO) fine-tuning for critical parameters with zeroth-order (ZO) corrections for less sensitive layers, as inspired by Hi-ZFO (Jin et al., 2026).
   - ZO optimization introduces beneficial stochasticity to escape local minima.
2. **Task Prototype-Based Retrieval**:
   - Utilizes task prototypes, as proposed by Oh et al. (2026), to embed task-specific characteristics and quantify associations between tasks.
   - Prototypes are refined using hierarchical embedding models.
3. **Entropy-Guided Alignment**:
   - Adapts Khoury et al. (2026)’s entropy-guided prompt weighting to adjust predictions dynamically, ensuring robust task alignment even in low-resource scenarios.

#### Experimental Setup
We evaluated HAFTAR on three benchmarks:
1. **ReasonTabQA** (Pan et al., 2026): A complex TableQA benchmark with multi-table structures.
2. **CLaS-Bench** (Gurgurov et al., 2026): A multilingual benchmark for evaluating language steering.
3. **Synthetic Reasoning Tasks**: Derived from the GanitLLM corpus (Dipta et al., 2026), featuring difficulty-aware Bengali mathematical reasoning.

#### Metrics
- **Accuracy**: Task-specific prediction accuracy.
- **Semantic Relevance**: Measured via harmonic-mean steering scores (Gurgurov et al., 2026).
- **Efficiency**: Computational overhead and convergence time.

---

### Results
| **Benchmark**       | **Baseline Accuracy** | **HAFTAR Accuracy** | **Semantic Relevance** | **Efficiency Gain** |
|----------------------|------------------------|----------------------|-------------------------|----------------------|
| ReasonTabQA          | 47.2%                 | **58.9%**           | 0.67                   | **1.5x**            |
| CLaS-Bench           | 65.4%                 | **75.3%**           | **0.81**               | 1.2x                |
| Synthetic Reasoning  | 72.1%                 | **81.4%**           | 0.78                   | **2.0x**            |

---

### Discussion
HAFTAR demonstrated robust performance across all benchmarks, achieving significant improvements in task accuracy, semantic relevance, and computational efficiency. The hierarchical optimization framework effectively mitigated issues of local minima and instability in fine-tuning, while entropy-guided alignment ensured robust task adaptation in low-resource settings. Notably, HAFTAR's integration of task prototypes enabled improved generalization across multilingual and cross-domain tasks.

The results underscore the importance of combining optimization techniques with retrieval-based approaches to address the inherent complexities of fine-tuning LLMs in real-world scenarios.

---

### Conclusion
This paper introduced HAFTAR, a novel framework for hierarchical optimization and task adaptation in LLMs. By integrating first- and zeroth-order optimization, task prototype-based retrieval, and entropy-guided alignment, HAFTAR achieved superior performance across diverse benchmarks. The findings highlight the potential of hybrid approaches to enhance fine-tuning, task retrieval, and generalization for LLMs, paving the way for more robust and efficient language model deployment in complex domains.

---

### Contributions
1. **Novel Framework**: Proposed HAFTAR, a hierarchical optimization framework for fine-tuning and adaptation in LLMs.
2. **Task Prototype Integration**: Leveraged task prototypes for robust multi-task learning under partial supervision.
3. **Entropy-Guided Alignment**: Adapted entropy-guided methods for improved task alignment and generalization.
4. **Comprehensive Evaluation**: Demonstrated the efficacy of HAFTAR across ReasonTabQA, CLaS-Bench, and synthetic reasoning benchmarks.

--- 

